% History:
% 30Dec91	created							dak
% 3 Jan92	complete						dak

\baselineskip 18truept

\centerline {\bf Transputer Common Library Routines}
\centerline {\bf 30 Dec 91}

\vskip 20truept
\item {\bf Introduction}
	These are the routines and macros that make up the ``common'' 
	subdirectory.  To use them, put a {\tt include ../com.make} in
	your makefile and tack {\tt \$(COMMON\_OBJS) \$(OPTIONAL\_OBJS)} onto
	the end of your {\tt OBJS} line.

\vskip 20truept
\item {\bf commands.c}  
	This is a file that you must put in your own directory.  It has
	declarations for all of the variables and commands that the libraries
	support.  This file should include {\tt <conc.h>} and {\tt "common.h"}
	in that order.  These are the things that must be defined in {\tt
	commands.c}:
\itemitem {\bf MEMORY}  
		This tells the memory management routines how much memory is
		available.  This is just a macro which defines some variables,
		and it wants the number of bytes as an argument.  An example
		is {\tt MEMORY(3*1024*1024 + 512*1024);} to reserve 3.5
		megabytes of memory.
\itemitem {\bf command}
		Is an array of type {\tt command\_form} which contains pointers
		to functions.  These functions correspond to the protocol
		commands, one function for each command.  If a command
		is to be ignored you can use a {\tt NULL} function pointer
		to have the controller dispose of te message.  The format of
		the functions to be called is discussed in the {\bf control.c}
		section.
\itemitem {\bf num\_commands}
		Tells the controller how many positions there are in the command
		array.  This is normally 18 for all messages.
\itemitem {\bf var\_ptr}
		Is an array of integer pointers which correspond to the
		protocol variables.
\itemitem {\bf setv\_cmnds}
		Is an array of functions pointers (type {\tt command\_form})
		which match up with the variables in the {\tt var\_ptr} array.
		These functions are called when a variable is being set.
		These can all be set to {\tt set\_var} which is described in
		the {\bf var\_sr.c} section.
\itemitem {\bf readv\_cmnds}
		Is an array of functions pointers (see above) for actions to
		take when reading a variable.  These can be set to 
		{\tt read\_var}.
\itemitem {\bf chan\_list}
		Is an array of channels for the controller to watch.  Because
		you will probably use dynamically allocated channels, you will
		have to put {\tt NULL} placeholders in and fill it in at 
		runtime.

\vskip 20truept
\item {\bf talk.c}
	contains {\tt reader} and {\tt writer} functions that must be run as
	processes.
\itemitem {\bf reader}
		watches all input channels defined in the {\tt chan\_list}
		array.
\itemitem {\bf writer}
		watches the {\tt Reader\_to\_Writer} channel for things to
		send.  It calls the router function to decide which channel
		to send through.  If the message is sent to the current
		transputer, it will be sent through the 
		{\tt Writer\_to\_Control} channel for the controller to deal
		with (see {\bf control.c}).

\vskip 20truept
\item {\bf control.c}
	handles all messages.  The {\tt main()} program should call 
	{\tt control()} as the last thing that it does.  Control never
	returns, rather it receives messages through the 
	{\tt Writer\_to\_Control} channel and dispatches the according to the
	{\tt command} structure defined in {\bf commands.c}.  If the function
	pointer has a {\tt NULL} value then the message will be discarded.
	Control also disposes of the buffer when the function returns.  Each
	function that is called should receive two arguments, {\tt int header},
	and {\tt int *buf} in that order.  The header can be decoded using the
	routines in {\tt protocol.h} if necessary.

\vskip 20truept
\item {\bf var.c}
	deals with variable set and reads.  It call the appropriate routine
	based on the {\tt readv\_cmnds} and {\tt setv\_cmnds} arrays defined
	in {\bf commands.c}.  The {\tt READ\_VAR} and {\tt SET\_VAR} commands
	should call {\tt var()}.

\vskip 20truept
\item {\bf var\_sr.c}
	has two routines for setting and reading of normal variables.  The 
	{\tt readv\_cmnds} and {\tt setv\_cmnds} arrays for normal variables
	should have pointers
	to {\tt read\_var} and {\tt set\_var()} respectively.

\vskip 20truept
\item {\bf mem.c}
	is a memory management package.  It uses the memory set aside in the
	{\bf commands.c} files using {\tt MEMORY}.  This is used to control
	where memory is allocated and to also provide a simple garbage
	collection mechanism.
\itemitem {\bf mem\_init}
		sets all of memory to be free.  This should be called from
		{\tt main()}.
\itemitem {\bf mem\_alloc}
		returns a pointer to a block of memory of the given size.  This
		routine returns a NULL if the memory cannot be allocated.
\itemitem {\bf mem\_free}
		returns a block of memory to the system heap.  If the memory is
		outside of the heap or was not allocated the routine returns
		without taking any action.

\vskip 20truept
\item {\bf common.h}
	provides function prototypes for all of the common routines.  This
	file also provides some useful macros and constants that configure the
	common routines.  One useful macro is {\tt SEND(header, buf, channel)}.

\vskip 20truept
\item {\bf monitor.c}
	has a watchdog process. ({\it not implemented})

\vskip 20truept
\item {\bf proc.c}
	Defines functions for downloading code on the fly.  The {\tt common.h}
	file also defines some macros that let your controller communicate with
	the downloaded process.  Specifically there is a control register,
	{\tt register}, and a channel, {\tt PROC\_to\_Control}.
\itemitem {\bf prog}  
	is the downloaded program.  This is really just an integer array of
	size {\tt PROG\_LEN} ints.  Just include this in the {\tt var\_ptr}
	array so that programs can be downloaded.
\itemitem {\bf kill\_proc}
	will ``kill'' the current process.  This relies on the process checking
	the contents of {\tt register} frequently.  When {\tt register} contains
	the value {\tt DIE}, it should send an int through the 
	{\tt PROC\_to\_Control} channel to signify completion and finally
	call ProcHalt().
\itemitem {\bf execute\_proc}
	starts the downloaded process.  It allocates a stack of size 
	{\tt P\_STACK}, sets {\tt register} to 0, and runs the process
	at low priority.
\itemitem {\bf register}
	is an int located in a static location that can be used to communicate
	with the downloaded process.  The process must recognize the {\tt DIE}
	command.
\itemitem {\bf PROC\_to\_Control}
	is a channel located in a static location and it may be used to
	communicate with the downloaded process.  The process should send
	an int out of this channel to signal that it is dying.

\vskip 20truept
\item {\bf queue.c}
	provides queue manipulation routines.  For sake of efficiency the
	queue routines do not provide error checking.  If this is necessary,
	you must make a call to emptyq yourself.
\itemitem {\bf q\_alloc}
		allocates an empty queue.  The queue has size {\tt QUEUE\_SIZE}.
\itemitem {\bf enq}
		adds an element (an integer size piece of data) to the tail of
		the queue.
\itemitem {\bf deq}
		removes and returns an element from the head of the queue.
\itemitem {\bf emptyq}
		returns a true value if the queue is empty, false if it is not.

\vskip 20truept
\item {\bf debug.c}
	provides a way of sending debugging messages up to the sun.
	Simply make a call to {\tt send\_debug(char *format, ...)} and it will
	do the rest.  The {\bf reader} must be watching {\tt DEBUG\_CHAN}.
\end
